% Autor: Dominik Harmim <xharmi00@stud.fit.vutbr.cz>


\documentclass[a4paper, 11pt]{article}


\usepackage[czech]{babel}
\usepackage[utf8]{inputenc}
\usepackage[left=2cm, top=2cm, text={17cm, 25cm}]{geometry}
\usepackage{times}
\usepackage{graphicx}
\usepackage[unicode, hypertexnames=false]{hyperref}


\begin{document}
	%%%%%%%%%%%%%%%%%%%%%%%%%%%%%%%% Titulek %%%%%%%%%%%%%%%%%%%%%%%%%%%%%%%%%%%
			\begin{center}
				{\Large
					Vysoké učení technické v~Brně \\
					Fakulta informačních technologií \\
				}
				{\includegraphics[width=0.4 \linewidth]{inc/FIT_logo.pdf}} \\

				{\LARGE
					Mikroprocesorové a~vestavěné systémy \\
					Projekt\,--\,Spínání světla dle intenzity \\[0.4cm]
				}

				{\large
					Jiří Prokop (xproko47) \\
					\texttt{xproko47@stud.fit.vutbr.cz} \\
					\today
				}
			\end{center}
	

	%%%%%%%%%%%%%%%%%%%%%%%%%%%%%%%% Úvod %%%%%%%%%%%%%%%%%%%%%%%%%%%%%%%%%%%%%%
	\section{Úvod}

        Cílem tohoto projektu bylo měřit úroveň svitu a podle toho nastavovat úroveň svitu 
        diody. Aplikace měla být uživatelsky nastavitelná. K přenosu hodnot byl určen protokol 
        MQTT.
        \newline
        Video s představením funkcionality lze nalézt zde: \url{https://drive.google.com/file/d/1RucK6XKDmvDGQYvzlt_hCG0MdMJrHIEZ/view?usp=sharing}


	%%%%%%%%%%%%%%%%%%%%%%%%%%%%%%%% Hardware %%%%%%%%%%%%%%%%%%%%%%%%%
        \section{HW součástky}
        \subsection{WeMos D2 R32 UNO ESP32}
            Deska integruje bezdrátový WiFi modul s čipem ESP-32 a zároveň obsahuje USB rozhraní pro napájení a 
            programování. Komunikaci zajišťuje řadič CH340. Samotný modul ESP-32 je vybaven 32bitovým 
            mikroprocesorem a nabízí mnoho GPIO pinů pro různé funkce dle naprogramovaného kódu.


        \subsection{BH1750}
            Snímač intenzity osvětlení BH1750 spolu s modulem senzoru osvětlení GY-302 určený pro mikrokontrolery.
            

        %%%%%%%%%%%%%%%%%%%%%%%%%%%%%%%% Zapojení hardware %%%%%%%%%%%%%%%%%%%%%%%%%
	\section{Příprava}


        \subsection{Zapojení}
	   \textbf{BH-1750} \\
            Sensor intenzity světla BH-1750 má 5 pinů z čehož stačilo zapojit pouze 4. Zapojil jsem je následovně:
            \begin{itemize}
                \item GND (ground) do pinu GND na esp32.
                \item SCL (clock) do SCL na esp
                \item SDA (data) do SDA na esp
                \item VCC (zdroj) do 3V3.
            \end{itemize}
            Zmíněným pátým nevyužitým pinem je ADDR, jehož zapojení záleží na použité konfiguraci senzoru.
            
          \vspace{6pt}
          \textbf{LED} \newline
          LED dioda je zapojena do analogového pinu IO2 a má připojený rezistor.
 
        \subsection{Vývojové prostředí}
        Při nastavování vývojového prostředí jsem následoval návod \cite{instalace_navod}, který byl k nalezení v
        IS VUT. Tímto bych za poskytnutí návodu chtěl poděkovat.

 
	%%%%%%%%%%%%%%%%%%%%%%%%%%%%%%%% Způsob řešení %%%%%%%%%%%%%%%%%%%%%%%%%
	\section{Způsob řešení}
            Projekt je implementován v jediném souboru main.c.
            Hlavní funkce programu je definována jako app\_main(). V této funkci jsou provedeny následující kroky:
            Inicializace LEDC pro ovládání LED.
            Inicializace Non-Volatile Storage (NVS) pro ukládání konfiguračních dat.
            Připojení k Wi-Fi síti pomocí funkce wifi\_connection(). Program čeká, dokud není získána IP adresa z 
            Wi-Fi.
            Načtení hraniční hodnoty pro LEDC z NVS nebo použití default hodnoty.
            Inicializace I2C pro komunikaci s periferiemi.
            Spuštění MQTT klienta pomocí funkce mqtt\_app\_start().
            Následně jsou v nekonečné smyčce odesílána sensorová data přes MQTT. Program získává hodnotu intenzity 
            světla ze senzoru pomocí funkce get\_data\_from\_sensor(). Pokud je hodnota -1, nastala chyba a čeká se na další měření. 
            Jinak je hodnota konvertována na řetězec a odeslána přes MQTT klienta. Intenzita světla je pak také použita pro nastavení jasu LED.
            Nakonec je ještě krátké čekání a smyčka se opakuje.
            

        \subsection{Konfigurace BH1750}
            Při konfiguraci jsem se řídil manuálem k této součástce \cite{sensor_svetla}. Jako mód jsem vybral
            'Continuously H-Resolution Mode', který zde byl doporučen. Při této konfiguraci je možno ze sensoru 
            snímat nové hodnoty jednou za 120ms. V programu jsem se ale rozhodl snímat je pouze jednou za 500ms,
            z důvodu zpoždění zasílaných MQTT zpráv. Při snímání každých 120ms se totiž často stávalo, že se
            zpráva nestihla odeslat včas, takže se jich následně odeslalo několik zároveň.

            V implementaci ostatních konfiguračních parametrů jsem se řídil podle přídladu uvedeného v dokumentaci
            ESP-IDF \cite{sensor_setup_example}.

        \subsection{Nastavení diody}
            Při nastavování pinů diody a při testování správnosti jejího zapojení jsem se inspiroval v 
            dokumentaci ESP-IDF \cite{led_blink}. Při následném nastavování svítivosti zase v programovacím
            návodu ESP-IDF \cite{led_control}. 
            % TODO

        \subsection{Připojení k WiFi a využití NVS}
            Pro připojení k síti WiFi jsem rovněž využil dokumentace ESP-IDF \cite{wifi}.
            Stejně tak pro nastudování, jak použít Non-Volatile Storage (NVS) k permanentnímu uložení dat 
            \cite{nvs_example}.

        \subsection{MQTT a uživatelské rozhraní}
            MQTT představuje jednoduchý protokol pro síťovou komunikaci typu 
            publish-subscribe mezi zařízeními (machine to machine) s funkcionalitou fronty zpráv. Je navržen pro 
            spojení s vzdálenými místy, kde existují zařízení s omezenými zdroji nebo omezenou šířkou pásma, jako 
            například v tomto projektu.

            Komunikaci přes MQTT jsem implementoval díky programovacímu návodu ESP-IDF \cite{mqtt}. Jako brokera jsem použil
            zdarma testovací službu EMQ \cite{mqtt_broker}.

            Jak již bylo dříve zmíněno, pomocí MQTT se odesílá každá naměřená hodnota. 
            Tu si přečte klient
            uživatelského rozhraní a zobrazí ji uživatelovi. Uživatel zde také může nastavit prahovou hodnotu pro
            spuštění diody. Po tom co ji esp obdrží ji pomocí NVS uloží do paměti a začne ji používat.
            
            Uživatelské rozhraní lze nalézt na této adrese: \url{https://www.stud.fit.vutbr.cz/~xproko47/IMP/}.
            
	%%%%%%%%%%%%%%%%%%%%%%%%%%%%%%%% Citace %%%%%%%%%%%%%%%%%%%%%%%%%%%%%%%%%%%%
        \newpage
	\bibliographystyle{czechiso}
	\renewcommand{\refname}{Literatura}
	\bibliography{documentation}
\end{document}
